\section{Conclusion}
This work highlights the importance of understanding the domain, autorater, and rubric during rubric creation. First, adding examples and context significantly improves human-autorater agreement in addition to autorater self-agreement, but this is dependent on the domain and autorater.
Second, rubrics that reduce confirmation bias tend to provide significant improvement for human-autorater agreement. 
Third, rubric complexity and aggregation methods across holistic and analytic rubrics influence human-autorater agreement.
Lastly, higher human inter-rater agreement contributes positively to human-autorater agreement. Practitioners aiming to use autoraters should carefully curate human annotation data and design rubrics that appropriately address the differences across domains and autoraters. Future work should explore a wider range of domains, autoraters, and rubrics to develop more comprehensive recommendations.

